\chapter{Finite semidefinite representation}
\label{ch:theorem}

Consider two separated quantum systems, with two binary measurements
on each. Their behaviour is specified by the eight expectations
\[
 q=(a_0,a_1,b_0,b_1,E_{00},E_{01},E_{10},E_{11}),\qquad
 p(\alpha,\beta\mid x,y)=
 \frac{1+\alpha a_x+\beta b_y+\alpha\beta E_{xy}}4.
\]
Let $\mathcal Q_{222}$ contain all such behaviours obtained from
finite-dimensional quantum states and binary positive-operator-valued
measurements. Local dimensions are unrestricted.

\begin{theorem}\label{thm:main}
There are finite integers $N,m$ and real symmetric matrices
$A_0,A_1,\ldots,A_8,B_1,\ldots,B_m$ of order $N$ such that
\[
 \mathcal Q_{222}=\left\{q\in\mathbb R^8:\ \exists y\in\mathbb R^m,\quad
 A_0+\sum_{\nu=1}^8q_\nu A_\nu+\sum_{j=1}^my_jB_j\succeq0\right\}.
\]
\end{theorem}

Thus every Bell functional in this scenario can be optimized over
one fixed finite semidefinite representation. We prove existence by
certifying all supporting inequalities. The certificate must remain
valid when the inequality and the quantum realization are perturbed
independently. This independence is expressed by a tensor product.

\begin{theorem}[Separated sums of squares]\label{thm:slack}
Let $K_c,K_t$ be labelled real closed extensions of $\mathbb R$.
If $f=f_0+\sum_{\nu=1}^8f_\nu q_\nu$ is nonnegative on
$(\mathcal Q_{222})_{K_c}$ and $q\in(\mathcal Q_{222})_{K_t}$, then
\[
 [f,q]:=f_0\otimes1+\sum_{\nu=1}^8f_\nu\otimes q_\nu
       \in\Sigma(K_c\otimes_{\mathbb R}K_t)^2.
\]
Here $\Sigma T^2$ denotes finite sums of squares in $T$.
The two fields may be labelled copies of the same extension.
\end{theorem}

The proof proceeds from the physical model to the supporting
inequalities, and then to their singular contacts.

Chapter~\ref{ch:model} first reduces the quantum set to
$\operatorname{conv}\Phi(S^1\times S^1\times S^3)$, where
$\Phi$ is the polynomial expectation map for two real qubits.
At most nine such points are needed in a mixture. This gives a
compact convex semialgebraic set and a finite formula for its
extension to any real closed field.

Chapter~\ref{ch:geometry} then reduces a valid inequality to
nonnegative combinations of events and canonical supporting
inequalities. Two closed families have explicit small matrix
representations. Every other extreme inequality has a stationary
zero and, after relabelling, one of two event supports:
\[
\begin{array}{ll}
\mathsf C_1:&\lambda=n=0,\quad d,m,C,A,\alpha>0,
                       \quad D,\beta\ge0,\quad D+\beta>0,\\
\mathsf C_2:&\lambda=\beta=0,\quad d,m,C,A,D,\alpha>0,
                       \quad n\ge0.
\end{array}
\]
Chapter~\ref{ch:entry} derives the signs and rational coordinates
of these contacts from the original stationary equations.

It remains to certify the corresponding quadratic form in the
state amplitudes. Chapter~\ref{ch:tensor} provides the common
algebraic step: a matrix with bounded tensor coefficients and
positive-definite real residue gives an exact sum of squares.
The local calculations express the full quadratic form in roots
for which this criterion applies. Chapter~\ref{ch:cover} gives
an exhaustive source partition; the subsequent proofs establish
its identities at every boundary. Vanishing event weights and
commuting measurement limits remain part of the argument.

Finally, finite mixtures preserve the certificates. If
$f=\sum_j a_jf_j$, $q=\sum_iw_iq_i$, $a_j,w_i\ge0$ and
$\sum_iw_i=1$, and if $[f_j,q_i]=\sum_k h_{ijk}^2$, then
\[
 [f,q]=\sum_{i,j,k}
       \bigl[(\sqrt{a_j}\otimes\sqrt{w_i})h_{ijk}\bigr]^2.
 \tag{M2}
\]
The normalization of the testing weights preserves the constant
term of $f$. Theorem~\ref{thm:slack} follows. Applying
Scheiderer's tensor-evaluation criterion
\cite[Corollary 3.19]{scheiderer} to two labelled copies of each
real closed field gives Theorem~\ref{thm:main}.

Throughout, source coefficients and contact coordinates belong
to $K_c$; testing measurements and states belong to $K_t$.
The symbols $\mathrm I,\mathrm U,\mathrm O$ distinguish positive
infinitesimals, positive bounded elements with bounded inverse,
and positive unbounded elements. Formula tags are local to each
section. Every sum and algebraic operation in a certificate is finite.
